% This is a Plain TeX document, NOT a LaTeX document.
% There is no preamble, no \documentclass, and no \begin{document}.

% 1. Setting Up Margins (Plain TeX style)
\hsize=6.5in
\vsize=9.0in
\hoffset=0.0in
\voffset=0.0in

% 2. Title Section
\font\titlefont=cmr12 scaled \magstep3
\font\boldfont=cmbx10 scaled \magstep1
\centerline{\titlefont Plain TeX Demonstration}
\vskip 0.5in

% 3. First Paragraph
This is a simple document typeset using standard {\boldfont Plain TeX} rather 
than LaTeX. In Plain TeX, paragraph breaks are created by leaving a blank line 
in the source file, just like in LaTeX.

\vskip 0.25in

% 4. Headings (Created manually using primitives)
{\boldfont 1. Mathematical Formulas}
\medskip
You can write inline math like $E = mc^2$, or displayed math equations using 
double dollar signs:
$$ \sum_{n=1}^{\infty} {1 \over n^2} = {\pi^2 \over 6} $$
Notice that instead of \LaTeX's \vbox{\tt\string\frac\{1\}\{n\^2\}}, Plain TeX uses the 
native {\tt\string\over} primitive.

\vskip 0.25in

% 5. Creating a List
{\boldfont 2. Bulleted Lists}
\medskip
Plain TeX uses the {\tt\string\item} macro for lists. It does not require wrapping 
them inside an environment:

\item{$\bullet$} This is the first item in a native TeX list.
\item{$\bullet$} This is the second item.
\item{$\bullet$} Notice there is no closing tag needed for this list.

\vskip 0.5in
\centerline{\it End of Document.}

% 6. The standard primitive to tell TeX to stop processing and output the page.
\bye
